% paper/sections/09d_defect_correction.tex
%
% §9.5  欠陥補正法（Defect Correction Method）の PPE への適用
%   — §9.2 で構築した CCD 離散演算子 L_H を欠陥補正で扱う枠組み
%   — Thomas 三重対角係数の導出：付録 \ref{app:sweep_thomas_coeff}
%   — 収束理論・スペクトル解析：付録 \ref{app:dc_convergence_theory}

\subsection{PPE への欠陥補正法の適用}
\label{sec:defect_correction_ppe}

\noindent\textbf{本節の位置づけ：}
§\ref{sec:defect_correction_main} で定義した欠陥補正法（DC）を，
PPE の CCD 高次演算子 $L_H$ と有限差分低次演算子 $L_L$ に具体化する．
本稿では smoothed-Heaviside 一括 PPE（連続圧力形式の PPE 経路）
および分相 PPE 解法（§~\ref{sec:split_ppe_framework}）の双方に形式適用するが，
両者の主張範囲は同じではない．
分相 PPE では各相の PPE が定数密度のラプラシアンとなるため
$L_H$/$L_L$ のスペクトル不整合が最小化され，
高次残差条件 $\eta_{\mathrm{DC}}\le\varepsilon_{\mathrm{PPE}}$ で受理することで
漸近精度 $\Ord{h^6}$ を得る（§~\ref{sec:ppe_dc_accuracy}；
境界条件種別の内訳は §\ref{sec:dc_pure_fccd_convergence} の表で
$\Ord{h^7}$ Dirichlet 境界条件 / $\Ord{h^5}$ Neumann 境界条件 / IIM 行補正で $\Ord{h^4}$ 以上 と区別）；
成分検証では少数回の補正で残差受理条件へ到達することを確認するが，
回数そのものは受理条件ではない．
smoothed-Heaviside 一括 PPE では $\rho$ 滑化幅
$\varepsilon_H$ が DC 収束率を律速し，
補正の回数は設計精度を保証する主受理条件ではなく，
低・中密度比の対照経路でも残差監視付き補助手順として扱う．
高密度比で残差が停滞する場合は，反復回数を増やす対象ではなく，
分相 PPE への経路選択診断である．
さらに実計算の PPE 射影では，固定緩和係数による更新が
高次残差を増やす場合，その補正は採用しない．
DC は「補正を何回加えたか」ではなく，
作用中の PPE 演算子に対する残差を減らす補正だけを受け入れる
射影精度の契約として用いる．

共通記法 \eqref{eq:dc_three_step} において，
PPE では $A_H=L_H$，$A_L=L_L$，$x=p$ とおく．
すなわち，§~\ref{sec:ccd_poisson_matrix} で構築した CCD 離散演算子 $L_H$（$\Ord{h^6}$）を
「\textbf{高精度で評価し，低精度で解く}」という原理で効率的に求解する
\textbf{欠陥補正法（Defect Correction Method）}~\cite{Stetter1978} を示す：
残差（欠陥）の評価には $L_H$ を用いて $\Ord{h^6}$ 精度を確保し，
線形系の求解には低次演算子 $L_L$（2次精度 有限差分）を用いる．
2次元の $L_L$ はブロック三重対角の疎行列であり，
基準解析では疎行列直接法で第2段の解法誤差を抑える．
一方，非一様格子では，
同じ $L_L$ 補正写像を複数の補正右辺に適用し，
第2段の直接法保証を保つ．
したがって本章の基準モデルは固定低次補正写像を用いる直接法であり，
第2段の線形解法誤差を DC 収束評価から切り離す．

§~\ref{sec:ccd_poisson_matrix} で CCD 空間離散化スキームを選択済みであり，
本節はその CCD 系を解く\textbf{線形求解法}の選択に特化する（収束特性の定量検証は §~\ref{sec:dc_iteration_accuracy}）．

\noindent\textbf{なぜ $L_H$ を直接解かないのか：}
$L_H$ は CCD ブロック Thomas 法に基づく密結合演算子であり，
$x$-$y$ 方向の結合が強いため ADI 的な方向分離が成立しない．
大域疎行列の陽的構成とその直接解法は付録~\ref{app:ccd_lu_direct} に示す．

\subsubsection{PPE 用欠陥補正の構成}
\label{sec:ppe_dc_algorithm}

離散化された PPE $L_H\,p = b$ に対し，
初期推定 $p^{(0)} = 0$ から以下の3段を反復する：
\begin{equation}
  \boxed{%
  \begin{aligned}
    &\textbf{第1段（欠陥の算出）：}\quad
      d^{(k)} = b - L_H\,p^{(k)}
    \\[4pt]
    &\textbf{第2段（補正の求解）：}\quad
      L_L\,\delta p^{(k+1)} = d^{(k)}
    \\[4pt]
    &\textbf{第3段（解の更新）：}\quad
      p^{(k+1)} = p^{(k)} + \alpha_k\,\delta p^{(k+1)}
  \end{aligned}
  }
  \label{eq:ppe_dc_three_step}
\end{equation}

各記号の定義：
\begin{itemize}
  \item $L_H$：CCD 法で離散化した可変係数ラプラシアン（$\Ord{h^6}$，§~\ref{sec:ccd_poisson_matrix}）．
        高次残差評価を担う．
  \item $L_L$：2次精度中心差分（有限差分）で離散化した近似ラプラシアン．
        2次元ではブロック三重対角の疎行列となる．
        基準解析ではこの疎行列を直接法で解き，第2段の低次補正方程式とする．
  \item $d^{(k)}$：欠陥（defect）— 現在の近似解 $p^{(k)}$ を
        高精度演算子 $L_H$ で評価した残差．
  \item $\delta p^{(k+1)}$：補正量（correction）— 低次演算子 $L_L$ の逆作用で得る．
  \item $\alpha_k$：補正方向に沿う更新係数．
        理想化した固定係数スペクトル解析では
        $\omega < 2/r_{\max} = 0.833$ が十分条件となる
        （付録~\ref{app:dc_convergence_theory}）が，
        実際の PPE 射影では固定 $\omega$ そのものを受理条件とはせず，
        下の残差最小化条件で各補正を判定する．
\end{itemize}

実装上は，補正方向 $\delta p$ に対して
\begin{equation}
  \alpha_\ast
  =
  \frac{\langle r^{(k)}, L_H\delta p^{(k+1)}\rangle_2}
       {\langle L_H\delta p^{(k+1)}, L_H\delta p^{(k+1)}\rangle_2},
  \qquad
  r^{(k)}=b-L_Hp^{(k)}
  \label{eq:ppe_dc_residual_minimizing_alpha}
\end{equation}
を第一候補にする．
これは $\|r^{(k)}-\alpha L_H\delta p^{(k+1)}\|_2^2$ を
その補正方向上で最小にする係数である．
$\alpha_\ast>0$ で残差が低下する場合に採用し，
そうでない場合は固定緩和候補の後退探索を行う．
どの候補でも $L_H$ 残差が下がらないときは，
その DC 小段は失敗として閉じ，残差を増やす圧力更新を速度補正へ渡さない．

\subsubsection{第2段の求解}
\label{sec:dc_sweep}

第2段の補正方程式 $L_L\,\delta p = d$ は，
$L_L$（smoothed-Heaviside 一括 PPE では変密度有限差分ラプラシアン，分相 PPE では相内定数密度有限差分ラプラシアン；
Neumann 境界条件とゲージ拘束付き）を
\textbf{低次有限差分補正求解法}で求解する：
\begin{equation}
  \delta p^{(k+1)} \approx L_L^{-1}\,d^{(k)}
  \label{eq:dc_step2_lu}
\end{equation}
ここで $L_L$ は 2 次元 Poisson のブロック三重対角疎行列であり，
単一方向の Thomas スイープだけで厳密に反転できる三重対角行列ではない．
行列形の基準解析ではゲージ拘束後の疎行列に対して直接解法を用いる~\cite{Kelley2003}．
一方，非一様格子では同じ $L_L$ を保存形有限差分フラックス形で組み立て，
初期補正と欠陥補正の複数右辺に同じ低次補正写像を適用する．
これにより第2段の直接法残差保証を保ったまま，
DC 反復の収束先を高次評価演算子 $L_H$ に揃える．
$L_L$ の Neumann 境界条件はゴーストセル反射 $p_{-1}=p_1$（左壁）で離散化し，
$L_H$（CCD）と $L_L$（有限差分）が同一の物理的 Neumann 条件を共有する
（境界スタンシルの導出は付録~\ref{app:sweep_thomas_coeff}）．
直接法の採用は第2段の $L_L$ 補正方程式の線形解法誤差を機械精度まで下げ，
反復法における許容残差 $\varepsilon_\mathrm{tol}$ の系統誤差蓄積を原理的に排除するためである．
このため本稿では，外側の欠陥 $d^{(k)}=b-L_Hp^{(k)}$ を高次演算子で評価しつつ，
低次補正方程式を直接法で解く構成を基準とする．
方向分離（ADI 型）スウィープとの比較解析は §~\ref{sec:verify_dc_spectral} および付録~\ref{app:dc_convergence_theory} を参照．

\subsubsection{PPE における収束先の精度保証}
\label{sec:ppe_dc_accuracy}

欠陥補正法の本質的な性質は，
\textbf{左辺に低次演算子 $L_L$ を用いても，高精度演算子 $L_H$ の精度が維持される}点にある．
収束条件 $\|d^{(k)}\| < \varepsilon_\mathrm{tol}$ を満たした時点で，
\begin{equation}
  L_H\,p^{(k)} \approx b
  \quad\text{（精度 $\varepsilon_\mathrm{tol}$）}
\end{equation}
が成立し，最終圧力解は CCD 演算子 $L_H$ の設計精度に漸近する．
ここでいう精度保持は，$L_H$ に対する欠陥 $\|b-L_Hp^{(k)}\|$ が十分小さいまで
DC 反復を収束させる場合の漸近的な主張である．
$L_L$ は収束を駆動する「加速装置」の役割のみを担い，
収束先の精度は右辺の欠陥評価（$L_H$ による残差）が決定する．
したがって，補正回数そのものは精度受理条件ではない．
分相 PPE の相内 Poisson では残差低下で受理し，一括変密度 PPE の界面帯誤差や
残差停滞を補正回数で打ち消せるとは主張しない．
特に，高密度比・界面ジャンプを含む実計算では，
固定緩和で残差が増える補正を繰り返すと，
圧力補正が非物理的な面加速度として運動方程式へ入る．
したがって式~\eqref{eq:ppe_dc_residual_contract} の受理判定には，
各受理更新が残差を下げるという単調性条件を含める．

\paragraph{PPE 射影での残差受理条件．}
相内 Poisson の成分検証で有効な固定反復回数をそのまま
二相 圧力ジャンプ PPE の精度受理条件にすると，
外側 DC の高次残差が射影誤差として残り得る．
したがって受理条件は反復回数ではなく，
\begin{equation}
  \eta_{\mathrm{DC}}^{(k)}
  =
  \begin{cases}
    \|b-L_Hp^{(k)}\|_2/\|b\|_2, & \|b\|_2>0,\\
    \|b-L_Hp^{(k)}\|_2, & \|b\|_2=0
  \end{cases}
  \le \varepsilon_{\mathrm{PPE}}
  \label{eq:ppe_dc_residual_contract}
\end{equation}
である．$k_{\max}$ はこの不等式に到達するための上限であって，
それ自体が数値精度の証明ではない．
この受理条件により，圧力場の評価では
「楕円方程式を十分に解いていない」誤差と，
圧力ジャンプ 閉包・圧力代表・界面幾何に由来する誤差を分離して扱える．
すなわち式~\eqref{eq:ppe_dc_residual_contract} は，
PPE 解法の反復パラメータではなく，射影段階の数値精度を判定する
離散残差条件である．

この精度保証は反復回数 $k$ と密接に関連する：
$k=1$ では $L_L$ の $\Ord{h^2}$ が支配的だが，
$k$ を増やすごとに $L_H$ の高次精度が発現し，
\textbf{等密度・均一格子・周期境界条件の理想化}では $k \geq 3$ で $\Ord{h^6}$ に漸近する
（定量的検証は §~\ref{sec:dc_iteration_accuracy} に示す）．
\textbf{超収束（Dirichlet 境界条件）}：
Dirichlet 境界条件 + 等密度の理想化下では，
$L_H$ の補正方程式 $L_L\delta=b-L_H p^{(k)}$ で内点欠陥が
$L_H$ の打切り誤差より 1 次低い形で打ち消し合うため，
$k\geq 3$ で局所 $\Ord{h^7}$ の超収束が観測される
（詳細は §\ref{sec:verify_split_ppe_dc}）．
本稿の主漸近主張は $\Ord{h^6}$；$\Ord{h^7}$ は Dirichlet 限定の超収束効果として
別途明示する．
\textbf{非理想条件下の劣化：}
(i) GFM 接続付き分相 PPE では界面条件の $\Ord{h^3}$ が律速して
$\Ord{h^4}$--$\Ord{h^5}$ に劣化；
(ii) Neumann 境界条件の境界閉包（§\ref{sec:ccd_bc}）では一般に $\Ord{h^5}$；
(iii) 非一様格子（$\alpha > 1$）では局所 $\Ord{h^4}$ に律速される．
§\ref{sec:verify_dc_spectral} の数値表は理想化条件と非理想条件の両方を示す．
反復行列のスペクトル解析（付録~\ref{app:dc_convergence_theory}）は，
等密度・固定係数・理想境界条件における固定緩和の参照範囲を与える：
$\omega \in (0,\, 0.833)$ が線形収束の十分条件であり，
$\omega \in \{0.3, 0.5, 0.7\}$ は同条件で良好な収束を示す
（実験検証: §~\ref{sec:verify_dc_spectral}）．
ただしこれは式~\eqref{eq:ppe_dc_residual_minimizing_alpha} の
残差低下受理を置き換えるものではない．
高密度比・界面ジャンプ・非一様格子を含む PPE では，
固定緩和値は後退探索の候補であって，生産計算の採否基準ではない．

\phantomsection\label{sec:dc_comparison}% backward-compat
\phantomsection\label{sec:dc_pure_fccd_outer_inner}% backward-compat

% ─────────────────────────────────────────────────────────────────────────────
\subsubsection{PPE 構成別の収束率（DC 残差受理条件と \texorpdfstring{$k$}{k} 上界）}
\label{sec:dc_pure_fccd_convergence}
% ─────────────────────────────────────────────────────────────────────────────

\noindent
3 種の PPE 構成（節点 CCD PPE，smoothed-Heaviside 一括 PPE，
分相 PPE + DC + CCD/FCCD 勾配）に対する収束率は，
GFM 律速・IIM 行補正・境界条件による上界に従う：
\begin{center}
\small
\renewcommand{\arraystretch}{\PaperTableArrayStretch}
\begin{tabular}{@{}p{0.30\linewidth}p{0.28\linewidth}p{0.34\linewidth}@{}}
\hline
PPE 構成 & 主張範囲 & 律速要因 \\
\hline
節点 CCD PPE（滑らかな製造解） & $\Ord{h^6}$（Dirichlet では $\Ord{h^7}$ 超収束を観測） & 滑らかな製造解条件のみ \\
smoothed-Heaviside 一括 PPE & 高解像度側でプラトー化し得る & $\rho^{-1}$ 遷移層と積の微分項 \\
分相 PPE + DC + CCD/FCCD 勾配 & 相内部は $\Ord{h^6}$（Dirichlet 境界条件の超収束 $\Ord{h^7}$，Neumann 境界条件の上界 $\Ord{h^5}$） & GFM/HFE/曲率/IIM 行補正の界面律速 \\
\hline
\end{tabular}
\end{center}
ここで「分相 PPE + DC + CCD/FCCD 勾配」は §\ref{sec:split_ppe_framework} の
圧力ジャンプ型閉包を指し，純 FCCD 高次極限そのものと同義ではない．
純 FCCD 高次極限は §\ref{sec:pure_fccd_dns} で述べる参照構成であり，
本表の相内 PPE・HFE・Balanced--Force 基本構成を全演算子へ拡張する構成である．
